\scanpage{407}
\mspage{32\deleted{\unsure{0}}1}
% archive sheet number: 406

% 1–2
\deleted{On va \unsure{maintenant} distinguer les deux cas

2 A) On a $\alpha=\alpha_0+1$.}

% 3–5
On va d'abord s'occuper des deux premières compatibilités $1_\alpha,2_\alpha$, en termes de \unsure{la} donnée de
% 6
\[
i_{\alpha_0}\in\operatorname{Ob} I,\quad \unsure{j|J}_{\alpha_0}\in\operatorname{Ob} J',\quad u_{\alpha_0}:\varphi(i_{\alpha_0})\longrightarrow j_{\alpha_0}.
\]

% 7–10
\deleted{\unsure{C|L}}\unsure{L|C}es \inserted{deux} compatibilités \unsure{voulues} pour les deux \unsure{objets} $i_\alpha,j_{\deleted{\unsure{\beta}}\alpha}$ et les quatre flèches $\tau_{\alpha,\alpha_0},\tau'_{\alpha,\alpha_0},v_{\alpha_0},u_\alpha$ sont indiquées dans le diagramme
% 11–13
\diagram{\begin{tikzcd}[ampersand replacement=\&,column sep=small]
\varphi(i_{\alpha_0})\arrow[r,"u_{\alpha_0}"]\arrow[rr,bend right=27,"{\varphi(\underline{\tau_{\alpha,\alpha_0}})}"'] \& j_{\alpha_0}\arrow[r,"{\underline{v_{\alpha_0}}\,?}"]\arrow[rr,bend left=27,"{\tau'_{\alpha,\alpha_0}}"] \& \varphi(\underset{?}{\underline{i_\alpha}})\arrow[r,"{\underline{u_\alpha}\,?}"] \& \underset{?}{j_\alpha}\qquad(*)
\end{tikzcd}}

% 14–16
On \unsure{note|voit} que comme $\varphi$ est cofinal, $\exists\,i_{\unsure{\alpha|?}}=i_{\deleted{\unsure{\alpha}}}\in\operatorname{Ob} I$ et $j_{\alpha_0}\xrightarrow{v_{\alpha_0}}\varphi(i_{\deleted{\unsure{\alpha}}})$. Mais on veut que le composé
% 17
\diagram{\begin{tikzcd}[ampersand replacement=\&,column sep=small]
w:\quad\varphi(i_{\alpha_0})\arrow[r,"u_{\alpha_0}"] \& j_{\alpha_0}\arrow[r,"v_{\alpha_0}"] \& \varphi(i_{\deleted{\unsure{\alpha}}})
\end{tikzcd}}
% 18–21
soit de la forme $\varphi(\tau)$, où $\tau=\tau_{\alpha,\alpha_0}:i_{\alpha_0}\longrightarrow i_\alpha$ est une flèche dans $I$. Comme $\varphi$ est un foncteur cofinal de catégories filtrantes, on sait
